% ============================================================
\chapter{Repr\'esentations Textuelles Classiques}
\label{ch:representations}
% ============================================================

Un algorithme ne lit pas. Il calcule. Pour qu'un classifieur, un moteur de
recherche ou un syst\`eme de traduction puisse travailler sur du texte, il
faut d'abord transformer les mots en nombres --- ou plus pr\'ecis\'ement, en
\emph{vecteurs}. Comment effectuer cette transformation~? C'est l'une des
questions fondatrices du TAL, et les r\'eponses ont \'evolu\'e
consid\'erablement~: des repr\'esentations \og sac de mots\fg{} des ann\'ees
1990 aux plongements denses de Word2Vec, en passant par la pond\'eration
TF-IDF.

\begin{intuition}[title=\textbf{Du texte aux vecteurs}]
Les algorithmes d'apprentissage automatique op\`erent sur des vecteurs
num\'eriques, pas sur du texte brut. Ce chapitre pr\'esente les m\'ethodes
classiques pour transformer un document textuel en une repr\'esentation
vectorielle exploitable par un classifieur ou un moteur de recherche.
\end{intuition}

% ------------------------------------------------------------
\section{Repr\'esentation sac-de-mots}
\label{sec:rep:bow}
% ------------------------------------------------------------

\begin{definition}[Sac-de-mots (Bag of Words)]
Le mod\`ele \emph{sac-de-mots} repr\'esente un document $d$ comme un vecteur
$\mathbf{x}_d \in \N^V$ o\`u la $j$-\`eme composante est le nombre
d'occurrences du terme $t_j$ dans $d$ :
\[
x_{d,j} = \text{count}(t_j, d)
\]
\end{definition}

Ce mod\`ele ignore compl\`etement l'ordre des mots. La phrase \og le chat mange
la souris \fg et \og la souris mange le chat \fg ont la m\^eme
repr\'esentation.

\begin{remarque}
Malgr\'e sa simplicit\'e, le mod\`ele sac-de-mots reste utile comme ligne de
base (\emph{baseline}) et dans des applications o\`u l'ordre des mots est
moins important (classification th\'ematique, filtrage de documents).
\end{remarque}

% ------------------------------------------------------------
\section{Pond\'eration TF-IDF}
\label{sec:rep:tfidf}
% ------------------------------------------------------------

La fr\'equence brute ne refl\`ete pas l'importance discriminative d'un terme.
Les mots tr\`es fr\'equents (``le'', ``de'', ``et'') dominent la repr\'esentation
sans apporter d'information s\'emantique.

\begin{definition}[TF -- Term Frequency]
La \emph{fr\'equence du terme} $t$ dans le document $d$ est :
\[
\text{tf}(t, d) = \frac{\text{count}(t, d)}{\sum_{t' \in d} \text{count}(t', d)}
\]
\end{definition}

\begin{definition}[IDF -- Inverse Document Frequency]
La \emph{fr\'equence inverse de document} du terme $t$ dans un corpus
$\mathcal{D}$ de $N$ documents est :
\[
\text{idf}(t, \mathcal{D}) = \log \frac{N}{\abs{\{d \in \mathcal{D} : t \in d\}}}
\]
\end{definition}

\begin{definition}[TF-IDF]
La pond\'eration \emph{TF-IDF} combine les deux mesures :
\[
\text{tf-idf}(t, d, \mathcal{D}) = \text{tf}(t, d) \times \text{idf}(t, \mathcal{D})
\]
\end{definition}

\begin{theoreme}[Propri\'et\'es de TF-IDF]
Soit $\mathbf{v}_d \in \R^V$ le vecteur TF-IDF du document $d$. Alors :
\begin{enumerate}
    \item $\text{tf-idf}(t, d) = 0$ si $t \notin d$ ;
    \item $\text{tf-idf}(t, d)$ est \'elev\'e si $t$ est fr\'equent dans $d$
          mais rare dans le corpus ;
    \item $\text{idf}(t) = 0$ si $t$ appara\^it dans tous les documents.
\end{enumerate}
\end{theoreme}

\begin{formules}[title=\textbf{Variantes de TF-IDF}]
\begin{align}
\text{tf}_{\log}(t,d) &= 1 + \log(\text{count}(t,d)) & \text{(tf logarithmique)} \\
\text{idf}_{\text{smooth}}(t) &= \log\frac{N+1}{\text{df}(t)+1} + 1 & \text{(idf liss\'e)} \\
\text{tf-idf}_{\text{norm}} &= \frac{\text{tf-idf}(t,d)}{\norm{\mathbf{v}_d}} & \text{(normalis\'e $L_2$)}
\end{align}
\end{formules}

\begin{codeblock}[title=\textbf{TF-IDF avec scikit-learn}]
\begin{minted}{python}
from sklearn.feature_extraction.text import TfidfVectorizer

corpus = [
    "Le chat dort sur le tapis",
    "Le chien court dans le jardin",
    "Le chat et le chien jouent ensemble"
]

vectorizer = TfidfVectorizer()
X = vectorizer.fit_transform(corpus)
print(f"Forme de la matrice : {X.shape}")
print(f"Vocabulaire : {vectorizer.get_feature_names_out()}")
print(f"Matrice TF-IDF (dense) :\n{X.toarray().round(2)}")
\end{minted}
\end{codeblock}

% ------------------------------------------------------------
\section{Mod\`eles $n$-grammes}
\label{sec:rep:ngrams}
% ------------------------------------------------------------

\begin{definition}[$n$-gramme]
Un \emph{$n$-gramme} est une sous-s\'equence contigu\"e de $n$ tokens extraite
d'un texte. Pour une s\'equence $(x_1, \ldots, x_T)$, les $n$-grammes sont :
\[
\{(x_t, x_{t+1}, \ldots, x_{t+n-1}) : t = 1, \ldots, T - n + 1\}
\]
\end{definition}

\subsection{Mod\`ele de langue $n$-gramme}

Un mod\`ele de langue $n$-gramme approche la probabilit\'e conditionnelle par
une hypoth\`ese de Markov d'ordre $n-1$ :

\begin{definition}[Mod\`ele de langue $n$-gramme]
\[
P(x_t \mid x_1, \ldots, x_{t-1}) \approx P(x_t \mid x_{t-n+1}, \ldots, x_{t-1})
\]
estim\'ee par maximum de vraisemblance :
\[
\hat{P}(x_t \mid x_{t-n+1}^{t-1}) = \frac{\text{count}(x_{t-n+1}, \ldots, x_t)}{\text{count}(x_{t-n+1}, \ldots, x_{t-1})}
\]
\end{definition}

\begin{attention}[title=\textbf{Probl\`eme des z\'eros}]
Si un $n$-gramme n'appara\^it jamais dans le corpus d'entra\^inement, sa
probabilit\'e estim\'ee est nulle, ce qui rend la perplexit\'e infinie. Les
techniques de lissage (Laplace, Kneser-Ney) sont indispensables.
\end{attention}

\subsection{Lissage de Laplace (add-$\alpha$)}

\begin{definition}[Lissage de Laplace]
Le lissage de Laplace ajoute un pseudo-comptage $\alpha > 0$ :
\[
\hat{P}_\alpha(x_t \mid x_{t-n+1}^{t-1}) = \frac{\text{count}(x_{t-n+1}^t) + \alpha}{\text{count}(x_{t-n+1}^{t-1}) + \alpha V}
\]
o\`u $V = \abs{\mathcal{V}}$.
\end{definition}

\subsection{Lissage de Kneser-Ney}

\begin{definition}[Lissage de Kneser-Ney]
Le lissage de Kneser-Ney utilise un discount absolu $\delta$ et une distribution
de continuation :
\[
P_{\text{KN}}(w \mid h) = \frac{\max(\text{count}(h, w) - \delta, 0)}{\text{count}(h)} + \lambda(h) \cdot P_{\text{cont}}(w)
\]
o\`u $P_{\text{cont}}(w) \propto \abs{\{h' : \text{count}(h', w) > 0\}}$
est la probabilit\'e de continuation et $\lambda(h)$ est un facteur de
normalisation.
\end{definition}

% ------------------------------------------------------------
\section{Similarit\'e entre documents}
\label{sec:rep:similarite}
% ------------------------------------------------------------

\begin{definition}[Similarit\'e cosinus]
La \emph{similarit\'e cosinus} entre deux vecteurs $\mathbf{u}, \mathbf{v} \in \R^V$ est :
\[
\cos(\mathbf{u}, \mathbf{v}) = \frac{\inner{\mathbf{u}}{\mathbf{v}}}{\norm{\mathbf{u}} \cdot \norm{\mathbf{v}}}
\]
\end{definition}

\begin{proposition}[Propri\'et\'es de la similarit\'e cosinus]
\begin{enumerate}
    \item $\cos(\mathbf{u}, \mathbf{v}) \in [-1, 1]$
    \item $\cos(\mathbf{u}, \mathbf{v}) = 1 \iff \mathbf{u} = \lambda \mathbf{v}$
          pour $\lambda > 0$
    \item Pour des vecteurs TF-IDF (composantes $\geq 0$) :
          $\cos(\mathbf{u}, \mathbf{v}) \in [0, 1]$
\end{enumerate}
\end{proposition}

\begin{definition}[Distance de Jaccard]
Pour deux ensembles de termes $A$ et $B$ :
\[
J(A, B) = \frac{\abs{A \cap B}}{\abs{A \cup B}}, \qquad d_J(A, B) = 1 - J(A, B)
\]
\end{definition}

% ------------------------------------------------------------
\section{Matrice documents-termes et SVD}
\label{sec:rep:svd}
% ------------------------------------------------------------

La matrice documents-termes $\mathbf{M} \in \R^{N \times V}$ est g\'en\'eralement
tr\`es creuse ($> 99\%$ de z\'eros). La d\'ecomposition en valeurs singuli\`eres
(SVD) permet de r\'eduire la dimensionnalit\'e.

\begin{definition}[Analyse s\'emantique latente (LSA)]
L'analyse s\'emantique latente (LSA) applique la SVD tronqu\'ee \`a la matrice
TF-IDF :
\[
\mathbf{M} \approx \mathbf{U}_k \boldsymbol{\Sigma}_k \mathbf{V}_k^\top
\]
o\`u $k \ll \min(N, V)$. Les documents sont alors repr\'esent\'es dans
$\R^k$ par les lignes de $\mathbf{U}_k \boldsymbol{\Sigma}_k$.
\end{definition}

\begin{theoreme}[Optimalit\'e de la SVD tronqu\'ee (Eckart-Young)]
La matrice de rang $k$ qui minimise l'erreur de Frobenius
$\norm{\mathbf{M} - \hat{\mathbf{M}}}_F$ est donn\'ee par la SVD tronqu\'ee
$\hat{\mathbf{M}} = \mathbf{U}_k \boldsymbol{\Sigma}_k \mathbf{V}_k^\top$.
\end{theoreme}

\begin{codeblock}[title=\textbf{LSA avec scikit-learn}]
\begin{minted}{python}
from sklearn.feature_extraction.text import TfidfVectorizer
from sklearn.decomposition import TruncatedSVD

vectorizer = TfidfVectorizer(max_features=10000)
X_tfidf = vectorizer.fit_transform(corpus)

svd = TruncatedSVD(n_components=100, random_state=42)
X_lsa = svd.fit_transform(X_tfidf)
print(f"Variance expliquee : {svd.explained_variance_ratio_.sum():.2%}")
print(f"Dimension reduite : {X_lsa.shape}")
\end{minted}
\end{codeblock}

% ------------------------------------------------------------
\section{Repr\'esentations par caract\`eres $n$-grammes}
\label{sec:rep:char}
% ------------------------------------------------------------

Les $n$-grammes de caract\`eres capturent la morphologie et sont robustes aux
fautes d'orthographe :

\begin{exemple}
Le mot \og bonjour \fg avec $n=3$ produit les trigrammes :
\texttt{\#bo}, \texttt{bon}, \texttt{onj}, \texttt{njo}, \texttt{jou},
\texttt{our}, \texttt{ur\#} (o\`u \texttt{\#} marque les fronti\`eres).
\end{exemple}

% ------------------------------------------------------------
\section{Mod\`ele BM25}
\label{sec:rep:bm25}
% ------------------------------------------------------------

\begin{definition}[BM25]
Le score BM25 du document $d$ pour la requ\^ete $q = \{q_1, \ldots, q_m\}$ est :
\[
\text{BM25}(d, q) = \sum_{i=1}^{m} \text{idf}(q_i) \cdot
\frac{\text{tf}(q_i, d) \cdot (k_1 + 1)}{\text{tf}(q_i, d) + k_1 \cdot \left(1 - b + b \cdot \frac{\abs{d}}{\text{avgdl}}\right)}
\]
o\`u $k_1$ et $b$ sont des hyperparam\`etres (typiquement $k_1 = 1.2$, $b = 0.75$)
et $\text{avgdl}$ est la longueur moyenne des documents.
\end{definition}

\begin{bonnepratique}[title=\textbf{BM25 en pratique}]
BM25 reste une m\'ethode de r\'ef\'erence pour la recherche d'information
(\emph{information retrieval}). M\^eme les syst\`emes RAG modernes utilisent
souvent BM25 comme premi\`ere \'etape de r\'ecup\'eration avant un reranking
neuronal.
\end{bonnepratique}

% ------------------------------------------------------------
\section{Limites des repr\'esentations creuses}
\label{sec:rep:limites}
% ------------------------------------------------------------

Les m\'ethodes pr\'esent\'ees dans ce chapitre souffrent de plusieurs limitations :

\begin{enumerate}
    \item \textbf{Haute dimensionnalit\'e} : $V$ peut atteindre $10^5$ \`a $10^6$.
    \item \textbf{Parcimonie} : les vecteurs sont tr\`es creux, rendant les
          mesures de distance peu fiables.
    \item \textbf{Absence de s\'emantique} : les mots synonymes ont des
          composantes orthogonales (``voiture'' $\perp$ ``automobile'').
    \item \textbf{Ind\'ependance contextuelle} : chaque mot a une repr\'esentation
          fixe ind\'ependante du contexte.
\end{enumerate}

Ces limitations motivent les plongements denses (chapitre~\ref{ch:plongements}),
qui repr\'esentent les mots dans un espace de faible dimension o\`u la proximit\'e
g\'eom\'etrique refl\`ete la similarit\'e s\'emantique.

% ------------------------------------------------------------
\section{Application : classification na\"ive de Bayes}
\label{sec:rep:naivebayes}
% ------------------------------------------------------------

\begin{definition}[Classifieur na\"if de Bayes multinomial]
Soit un document repr\'esent\'e par ses comptes de mots $\mathbf{x} = (x_1, \ldots, x_V)$.
Le classifieur na\"if de Bayes multinomial pr\'edit :
\[
\hat{y} = \argmax_{c \in \mathcal{Y}} \left[\log P(c) + \sum_{j=1}^{V} x_j \log P(t_j \mid c)\right]
\]
avec $P(t_j \mid c) = \frac{\text{count}(t_j, c) + \alpha}{\sum_{j'} \text{count}(t_{j'}, c) + \alpha V}$.
\end{definition}

\begin{codeblock}[title=\textbf{Classification na\"ive de Bayes}]
\begin{minted}{python}
from sklearn.feature_extraction.text import TfidfVectorizer
from sklearn.naive_bayes import MultinomialNB
from sklearn.pipeline import Pipeline
from sklearn.datasets import fetch_20newsgroups

data = fetch_20newsgroups(subset='train', categories=['sci.space', 'rec.sport.baseball'])
pipeline = Pipeline([
    ('tfidf', TfidfVectorizer(max_features=10000)),
    ('clf', MultinomialNB(alpha=0.1))
])
pipeline.fit(data.data, data.target)
print(f"Accuracy (train) : {pipeline.score(data.data, data.target):.2%}")
\end{minted}
\end{codeblock}

% ------------------------------------------------------------
\section{Exercices}
\label{sec:rep:exercices}
% ------------------------------------------------------------

\begin{exercice}[$\star$]
Calculez les vecteurs TF-IDF pour le corpus suivant :
$d_1$ = ``le chat dort'', $d_2$ = ``le chien dort'', $d_3$ = ``le chat et le chien''.
\end{exercice}

\begin{exercice}[$\star$]
Montrez que la similarit\'e cosinus est invariante par changement d'\'echelle
positif : $\cos(\lambda \mathbf{u}, \mu \mathbf{v}) = \cos(\mathbf{u}, \mathbf{v})$
pour $\lambda, \mu > 0$.
\end{exercice}

\begin{exercice}[$\star\star$]
Impl\'ementez un mod\`ele de langue bigramme avec lissage de Laplace. Calculez
la perplexit\'e sur un corpus de test.
\end{exercice}

\begin{exercice}[$\star\star$]
Montrez que le th\'eor\`eme d'Eckart-Young implique que la SVD tronqu\'ee donne
la meilleure approximation de rang $k$ au sens de la norme de Frobenius.
\end{exercice}

\begin{exercice}[$\star\star\star$]
Comparez les performances de TF-IDF + r\'egression logistique, TF-IDF + SVM,
et LSA + r\'egression logistique sur la classification de sentiments (dataset
IMDB). Analysez l'impact de la dimension $k$ de LSA.
\end{exercice}
